% !TEX program = xelatex
% À compiler sur Overleaf avec le compilateur XeLaTeX
% (Menu -> Compiler -> XeLaTeX) pour avoir la police Verdana.

\documentclass[12pt,a4paper]{article}

\usepackage{fontspec}
% Verdana si disponible, sinon une police sans-serif proche
\IfFontExistsTF{Verdana}{\setmainfont{Verdana}}{\setmainfont{DejaVu Sans}}

\usepackage[french]{babel}
\usepackage[margin=1.4cm]{geometry}
\usepackage{amsmath}
\usepackage{setspace}
\usepackage{enumitem}
\usepackage{array}

% Interligne 1,5
\onehalfspacing

% Espacement réduit autour des égalités successives
\setlength{\abovedisplayskip}{1pt}
\setlength{\belowdisplayskip}{1pt}
\setlength{\abovedisplayshortskip}{0pt}
\setlength{\belowdisplayshortskip}{0pt}

% ----- Pas de numérotation des pages -----
\pagestyle{empty}

\setlength{\parindent}{0pt}

% ----- Commandes personnalisées -----
\newcounter{exercice}
\newcommand{\exercice}{%
  \refstepcounter{exercice}%
  \vspace{0.4em}
  \noindent\underline{\textbf{Exercice \arabic{exercice}}}\par
  \vspace{0em}
}

\setlist[enumerate,1]{label=\alph*),itemsep=0pt,topsep=1pt}
\setlist[itemize,1]{leftmargin=1.6em,itemsep=0pt,topsep=1pt}

\begin{document}

% ----- En-tête -----
\begin{center}
  \underline{\Large\textbf{Correction du contrôle séquence 13 version 7}}\\[0.3em]
 
\end{center}

% =====================================================================
\vspace{3cm}\exercice
\begin{enumerate}
  \item \textbf{Oui.} $48  = 6 \times 8$ 
  \newline 6 est donc un diviseur de 48.
  \item \textbf{Non.} $184 \div 9 \approx 20{,}44$
  \newline le quotient n'est pas un entier,
        donc 184 n'est pas un multiple de 9.
  \item \textbf{Non.}  $21 < 63$, donc 21 n'est pas un multiple de 63.
  \item \textbf{Non.} $46 \div 4 = 11{,}5$ 
  \newline le quotient n'est pas un entier,
        donc 4 n'est pas un diviseur de 46.
\end{enumerate}

% =====================================================================
\vspace{1cm}\exercice
\begin{enumerate}
  \item \textbf{2 et 3} (acceptés aussi : 5 ; 7 ; 11 ; \dots).
        Un nombre premier a exactement deux diviseurs : 1 et lui-même.
  \item \textbf{4 et 6} (acceptés aussi : 1 ; 8 ; 9 ; \dots).
      
  \item Dans 16 ; 25 ; 5 ; 33 ; 65 ; 76 ; 89 ; 87, les deux nombres premiers sont
        \textbf{5} et \textbf{89}.
       
\end{enumerate}

% =====================================================================
\vspace{1cm}\exercice
\textbf{a) } $26 = 2 \times 13$, 
donc les diviseurs de 26 sont $\boxed{1\ ;\ 26\ ;\ 2\ ;\ 13\ }$.

\textbf{b) } $18 = 2 \times 3 \times 3$,
donc les diviseurs de 18 sont
$\boxed{1\ ;\ 18\ ;\ 2\ ;\ 3\ ;\ 6\ ;\ 9\ }$.

% =====================================================================
\vspace{1cm}\exercice


\noindent\begin{minipage}[t]{0.48\textwidth}
\textbf{a)}
\begin{align*}
35 &= 5 \times 7
\end{align*}

\textbf{b)}
\begin{align*}
28 &= 2 \times 14\\
   &= 2 \times 2 \times 7
\end{align*}
\end{minipage}\hfill
\begin{minipage}[t]{0.48\textwidth}
\textbf{c)}
\begin{align*}
260 &= 2 \times 130\\
    &= 2 \times 2 \times 65\\
    &= 2 \times 2 \times 5 \times 13\\
\end{align*}

\textbf{d)}
\begin{align*}
286 &= 2 \times 143\\
    &= 2 \times 11 \times 13
\end{align*}
\end{minipage}

% =====================================================================
\exercice
\noindent\begin{minipage}[t]{0.31\textwidth}
\textbf{a)} 
\begin{align*}
\dfrac{6}{21} &= \dfrac{3 \times 2}{3 \times 7}\\
              &= \dfrac{2}{7}
\end{align*}
\end{minipage}\hfill
\begin{minipage}[t]{0.31\textwidth}
\textbf{b)} 
\begin{align*}
\dfrac{24}{36} &= \dfrac{12 \times 2}{12 \times 3}\\
               &= \dfrac{2}{3}
\end{align*}
\end{minipage}\hfill
\begin{minipage}[t]{0.31\textwidth}
\textbf{c)} 
\begin{align*}
\dfrac{75}{195} &= \dfrac{15 \times 5}{15 \times 13}\\
                &= \dfrac{5}{13}
\end{align*}
\end{minipage}

% =====================================================================
\vspace{1cm}\exercice
Le nombre de paquets doit diviser à la fois 180 et 105 : c'est un diviseur commun à
180 et 105. Pour en faire un maximum, on cherche le plus grand diviseur commun à 180 et 105.

Avec les décompositions en facteurs premiers :
$180 = 2 \times 2 \times 3 \times 3 \times 5$ et $105 = 3 \times 5 \times 7$.
Le plus grand diviseur commun à 180 et 105 est le produit des facteurs premiers communs :
3 \times 5
                      = 15

Hulk et Batman peuvent donc faire au maximum \textbf{15 paquets},
chacun contenant  12 cookies et 7 tartelettes.

% =====================================================================
\vspace{1cm}\exercice
On utilise les critères de divisibilité :
\newline \textbf{par 2} : le chiffre des unités est 0, 2, 4, 6 ou 8 ;
\newline \textbf{par 3} : la somme des chiffres est divisible par 3
(sommes des chiffres : $29$ pour 845\,930 ; $72$ pour 888\,888\,888 ; $35$ pour 953\,585) ;
\newline \textbf{par 5} : le chiffre des unités est 0 ou 5 ;
\newline \textbf{par 10} : le chiffre des unités est 0.

\renewcommand{\arraystretch}{1.4}
\setlength{\tabcolsep}{2pt}
\begin{center}
\small
\begin{tabular}{|c|c|c|c|c|}
\hline
 & \textbf{Divisible par 2} & \textbf{Divisible par 3} & \textbf{Divisible par 5} & \textbf{Divisible par 10} \\
\hline
845\,930 & Oui & Non & Oui & Oui \\
\hline
888\,888\,888 & Oui & Oui & Non & Non \\
\hline
953\,585 & Non & Non & Oui & Non \\
\hline
\end{tabular}
\end{center}

\end{document}
